%% ============================================================
%%  Chapter 5: Conclusion
%% ============================================================
\chapter{Conclusion}
\label{chap:conclusion}

% GUIDANCE: 4--6 pages. Write this last.
% Structure: summary → contributions revisited → limitations →
%            future work.
% Be honest about limitations -- committees respect this far more
% than hand-waving.


% ── 5.1 ─────────────────────────────────────────────────────────
\section{Summary}
\label{sec:conclusion_summary}

% Briefly restate what you set out to do and what you found.
% Do NOT introduce new results here.

This thesis studied multi-robot coordination in warehouse environments under
communication constraints, using a hybrid allocation-and-planning stack that
combines LCBA with prioritized path planning and event-driven replanning.

The experimental results show that LCBA is competitive in stable, reliability-first
operating regimes and provides a robust decentralized baseline. Performance
degrades under overload or timeout-dominated conditions and is sensitive to map
geometry and load.

Steady-state experiments corroborate this pattern: LCBA maintains flat
throughput across the connectivity spectrum in stable regimes, while
stress-regime maps (Kiva, Kiva-Large, Shelf-Aisle) expose layout- and
congestion-induced failure modes (see Figures~\ref{fig:ss_kiva_heatmaps_stress},
\ref{fig:ss_kiva_large_heatmaps_stress}, and \ref{fig:ss_shelf_aisle_heatmaps_stress}).

The batch results support the main thesis claims of communication robustness,
preservation of near-optimality on connected graphs, and scalable behavior
under realistic loads. These findings are conditioned on map difficulty and
load; they do not imply unconditional superiority over all baselines in every
setting.


% ── 5.2 ─────────────────────────────────────────────────────────
\section{Contributions Revisited}
\label{sec:conclusion_contributions}

% Map each contribution listed in Section~\ref{sec:intro_contributions}
% to the specific evidence from Chapter 4 that supports it.

The experimental evidence supports the main contributions as follows: the
hybrid architecture is effective in the reliability-first batch regimes, the
dynamic replanning mechanism provides a meaningful responsiveness advantage
when communication and load remain within a usable range, and the overall
system behaves predictably enough to separate main-analysis maps from
stress-only maps. LCBA provides a practical, robust decentralized baseline for
warehouse coordination under communication constraints.


% ── 5.3 ─────────────────────────────────────────────────────────
\section{Limitations}
\label{sec:conclusion_limitations}

% Be honest. Committees will ask about limitations regardless.
% Naming them yourself shows rigour.

While the results demonstrate that LCBA is robust and practical in
stability-first operating regimes, several limitations should be acknowledged.

\textbf{Centralized path planning.}
Although task allocation is logically decentralized, the collision-avoidance
layer uses a shared space-time reservation table maintained by a central
orchestrator.
This constitutes a single point of failure and a communication bottleneck:
if the orchestrator is unavailable, path planning halts.
A fully decentralized deployment would require distributed reservation tables
or a per-component CBS-style planner, neither of which is implemented here.

\textbf{Simulation-only evaluation.}
All experiments are conducted in a discrete gridworld simulator.
Real warehouses differ in important respects: sensor noise, actuation
uncertainty, non-uniform robot speeds, and layout features (curved aisles,
varying shelf heights, human pedestrians) that the grid model does not capture.
The degree to which simulation results transfer to physical hardware
remains an open question.

\textbf{Generous SGA comparison on disconnected graphs.}
When the communication graph is disconnected, SGA is run independently on
each connected component.
This means SGA's sub-optimality guarantee applies per-component, not globally,
and the resulting SGA solution may itself be suboptimal relative to the true
global optimum.
Comparing LCBA against per-component SGA therefore represents a conservative
(generous) baseline for LCBA's sub-optimality ratio on disconnected graphs.

\textbf{Homogeneous fleet.}
All agents are assumed to have identical speeds, energy capacities, and task
capabilities.
Real fleets are often heterogeneous; the LCBA bid computation and energy
filter would need to be extended to handle agents with different operating
envelopes.

\textbf{Timestep ceiling at extreme conditions.}
At very high task loads combined with low communication ranges, the 3000-timestep
batch ceiling may terminate runs before the system has settled into a steady
operating regime.
This means that the worst-case completion data at $\comm = 4$,
$\text{tc} = 120$ reflects a timeout rather than a converged failure,
and caution is warranted when interpreting those data points.


% ── 5.4 ─────────────────────────────────────────────────────────
\section{Future Work}
\label{sec:conclusion_future}

Several directions would naturally extend this work.

\textbf{Fully decentralized collision avoidance.}
Replacing the central reservation table with a distributed planning scheme such as per-component CBS, a distributed RHCR variant or PIBT would remove the
single point of failure in the path planning layer and make the system fully
decentralized end-to-end.

\textbf{Time-windowed and deadline-constrained tasks.}
Real warehouse tasks often carry soft deadlines (orders must be processed
within a service-level window).
Extending the RPT utility function to penalise deadline violations would
enable LCBA to balance throughput against latency more explicitly, at the cost
of a more complex bid computation.

\textbf{Scaling to larger environments.}
The \texttt{shelf\_aisle} environment (200 agents, $161\times63$ grid)
represents the limit of what the current simulation stack can evaluate.
Experiments on even larger environments, or on real warehouse floor plans
extracted from operational data, would test whether the conclusions
transfer to production scale.

\textbf{Heterogeneous fleets.}
Extending the energy-budget filter and bid scoring to handle robots with
different speeds, battery capacities, and payload restrictions would
increase the practical applicability of LCBA in deployments where the
fleet is not uniform.

\textbf{Hardware validation.}
A hardware-in-the-loop experiment on a small fleet of physical robots would
validate the simulation findings and expose any gap between the grid model
and real robot dynamics (actuation noise, imprecise positioning, sensor
latency).

\textbf{Learned communication graph prediction.}
The current system reacts to topology changes after they occur.
A learning-based module that predicts imminent connectivity changes
from trajectory history could allow proactive reallocation before
agents disconnect, reducing the performance dip during topology transitions.
